\rhead{ML501: Natural Language Processing (NLP)}
\phantomsection 
\section*{Natural Language Processing (NLP) (ML501)}
\label{major:ML_5}

\noindent \small{\hyperref[main:scheme]{$\leftarrow$ Back to Scheme}} \vspace{0.5cm}

\noindent \textbf{Credits:} 3 (2-0-2)

\subsection*{Theory}
\noindent\textbf{Unit 1} \\
\textit{NLP Fundamentals and Text Preprocessing:} Language modeling and n-gram models, Regular expressions for tokenization, Sentence segmentation and normalization, Stemming, lemmatization, and part-of-speech tagging, Stopword removal and text normalization, Bag-of-words and TF-IDF representations, Character-level and subword tokenization (BPE, WordPiece), Unicode handling and multilingual text processing.

\vspace{0.3cm}

\noindent\textbf{Unit 2} \\
\textit{Word Embeddings and Sequence Models:} Word2Vec (skip-gram, CBOW), GloVe and fastText embeddings, Contextual embeddings (ELMo, Flair), RNN/LSTM/GRU for sequence modeling, Bidirectional encoders, Sequence labeling tasks (NER, POS tagging, chunking), CRF layer for structured prediction, Attention mechanisms (self-attention, multi-head attention).

\vspace{0.3cm}

\noindent\textbf{Unit 3} \\
\textit{Transformer Architecture and BERT:} Transformer model architecture (encoder-decoder, positional encoding), Self-attention and scaled dot-product attention, Multi-head attention and layer normalization, BERT pretraining objectives (MLM, NSP), RoBERTa, DistilBERT, and ALBERT variants, Fine-tuning strategies and domain adaptation, Sentence-BERT for semantic similarity.

\vspace{0.3cm}

\noindent\textbf{Unit 4} \\
\textit{Large Language Models and Prompt Engineering:} GPT architecture evolution (GPT-1 to GPT-4), Decoder-only transformers, In-context learning and few-shot prompting, Chain-of-thought reasoning, Retrieval-Augmented Generation (RAG) architecture, Knowledge graphs and dense retrieval, Prompt engineering techniques (zero-shot, few-shot, instruction tuning), Hallucination mitigation strategies.

\vspace{0.3cm}

\noindent\textbf{Unit 5} \\
\textit{Advanced NLP Systems and Evaluation:} Text generation evaluation (BLEU, ROUGE, BERTScore, human evaluation), Question answering systems (extractive, generative), Conversational AI (dialogue state tracking, response generation), Multilingual NLP (mBERT, XLM-R), Model deployment (TGI, vLLM), Ethical considerations (bias detection, toxicity classification, fairness evaluation), RAG evaluation frameworks.

\newpage

\subsection*{Practical}
\noindent\textbf{Lab 1: Text Preprocessing Pipeline} \\
Focuses on implementing complete text preprocessing with tokenization, normalization, and vectorization techniques.

\vspace{0.2cm}

\noindent\textbf{Lab 2: Word Embeddings and Similarity} \\
Focuses on training Word2Vec/GloVe, analogy tasks, and semantic similarity measurement.

\vspace{0.2cm}

\noindent\textbf{Lab 3: Sequence Models for NER/POS} \\
Focuses on BiLSTM+CRF implementation for named entity recognition and POS tagging.

\vspace{0.2cm}

\noindent\textbf{Lab 4: BERT Fine-tuning} \\
Focuses on fine-tuning BERT/RoBERTa for classification, NER, and question answering tasks.

\vspace{0.2cm}

\noindent\textbf{Lab 5: Transformer from Scratch} \\
Focuses on implementing multi-head self-attention and basic transformer encoder.

\vspace{0.2cm}

\noindent\textbf{Lab 6: LLM Prompt Engineering} \\
Focuses on zero-shot/few-shot prompting, chain-of-thought, and comparative evaluation.

\vspace{0.2cm}

\noindent\textbf{Lab 7: RAG Implementation} \\
Focuses on building complete RAG pipeline with dense retrieval and generation.

\vspace{0.2cm}

\noindent\textbf{Lab 8: Conversational AI Chatbot} \\
Focuses on dialogue systems with context tracking and response generation.

\vspace{0.2cm}

\noindent\textbf{Lab 9: Multilingual NLP Pipeline} \\
Focuses on cross-lingual transfer learning and zero-shot classification.

\vspace{0.2cm}

\noindent\textbf{Lab 10: Production NLP Application} \\
Focuses on end-to-end NLP system deployment with evaluation and monitoring.

\vspace{0.2cm}

\newpage
\rhead{Curriculum Schema}
